\section{Soundness}
\label{sec:soundness}

The rules of \Cref{sec:hoare} are lemmas about the semantics. Collecting them
into a syntactic proof system and proving that system sound is the only step
that is not a one-liner.

\begin{definition}[the proof system]
\label{def:derivable}
The judgement $\vdashH \triple{P}{c}{Q}$ is inductively generated by the five
structural rules --- the ones stated as \Cref{lem:hoare-skip},
\Cref{lem:hoare-assign}, \Cref{lem:hoare-seq}, \Cref{lem:hoare-if} and
\Cref{lem:hoare-while} --- read as introduction rules rather than as
statements about $\models$.
\end{definition}

\begin{theorem}[soundness]
\label{thm:soundness}
If $\vdashH \triple{P}{c}{Q}$ then $\triple{P}{c}{Q}$.
\end{theorem}

\begin{proof}
Induction on the derivation of \Cref{def:derivable}. Each case is the
corresponding lemma of \Cref{sec:hoare} applied to the induction hypotheses:
\Cref{lem:hoare-skip} and \Cref{lem:hoare-assign} are the base cases, and
\Cref{lem:hoare-seq}, \Cref{lem:hoare-if}, \Cref{lem:hoare-while} the step
cases. Nothing about \Cref{def:bigstep} is used directly here; it has all been
spent already.
\end{proof}

\begin{corollary}
\label{cor:derivable-holds}
If $\vdashH \triple{P}{c}{Q}$, the store $\sigma$ satisfies $P$, and
$\bigstep{c}{\sigma}{\tau}$, then $\tau$ satisfies $Q$.
\end{corollary}

\begin{proof}
Unfold \Cref{def:triple} in \Cref{thm:soundness}.
\end{proof}
